\section{Local weak bias and honest adaptation at all-valid data}
\label{sec:local-precision}

The broad-class obstruction in Section~\ref{sec:precision-boundary} leaves an
important question open: what length is necessary at favorable all-valid
data when coverage must also protect nearby weakly invalid configurations?
We answer this on an explicit Bernoulli linear submodel. The distinction
between source selection and honest inference follows the motivation of
\citet{guo2023}; moment matching also underlies adaptation costs in the
Gaussian model of \citet{wang2026}. The following calculation uses exact
patient-level Bernoulli likelihoods, with a trusted target and a deterministic
lower bound on the valid-source count.

\subsection{Define the local class before optimizing the method}

Fix $0<\rho<w<1$, $g=\lceil\rho K\rceil$ and a constant $M>0$. Consider
bounded linear arm means with at least $g$ valid sources per arm, and with
every source-assisted conditional mean bias bounded by
\[
 |\theta_{ja}-\mu_{t,X}^a|\le M r_{n,K},\qquad
 r_{n,K}=\min\{n_s^{-1/2}K^{-1/4},\ n_t^{-1/2}\}.
\]
There is no lower separation requirement. This defines the weak-bias subclass
for the present result; it does not constrain the main algorithm to know
which sources are invalid. For comparable target/source sample sizes, the
first term determines this local scale. In this class the arm dispersions
are at most $M^2r_{n,K}^2$.

\subsection{Exact two-world construction}

Let control means equal $.4$ at every site. Under $P_0$, all treated means,
including target, equal $\theta_0=.5$. Under the alternative, target and
valid-source treated means are $\theta_1=.5+\delta$, while invalid-source
treated means equal $.5-\lambda\delta$, where
\[
 \lambda=\frac{w}{1-w},\qquad \delta=c r_{n,K},\qquad
 0<c\le M(1-w).
\]
All site means are constant, hence exactly linear. Relative to the alternative
target, invalid-source bias is $-\delta/(1-w)$ and tends to zero. Draw each
site's valid flag with probability $w$, then condition on at least $g$ valid
flags. Every fixed configuration in this prior belongs to the declared local
class. Control means are valid everywhere. The target effects are $.1$ and
$.1+\delta$.

There are two experiments for which the likelihood calculation is explicit.
\begin{enumerate}
\item \emph{Conditional fixed arms:} each source contributes $n_s$ treated
Bernoulli observations and the target contributes $n_t$; put $\kappa=4$.
Independent control observations have identical laws in both worlds.
\item \emph{I.i.d. randomized patients:} each source has $n_s$ total patients
and the target has $n_t$. Covariates are i.i.d. from the same specified law in
the two worlds, $A\sim\operatorname{Bernoulli}(1/2)$ independently of $X$,
and outcomes have the constant means above. Put $\kappa=2$. Covariates are
ancillary for this likelihood comparison, and the population TATE equals
the constant effect. This is an unconditional random-design population
submodel, rather than a conditioning argument about one rare design.
\end{enumerate}

For one treated Bernoulli observation, the likelihood-ratio cross moment under
mean $.5$ is $1+4uv$ for shifts $u,v$. In the second experiment an untreated
observation has likelihood ratio one, so averaging over $A$ gives $1+2uv$.
Let $P_u^{(n_s)}$ denote one site's full observation law at treated shift $u$.
The unconditioned alternative site law is
$Q=wP_\delta^{(n_s)}+(1-w)P_{-\lambda\delta}^{(n_s)}$. Its chi-square
divergence from the all-valid null is exactly
\begin{align*}
 \chi_s^2
 &=\sum_{\ell=2}^{n_s}\binom{n_s}{\ell}
       (\kappa\delta^2)^\ell
       \{w+(1-w)(-\lambda)^\ell\}^2.\tag{L1}
\end{align*}
The degree-one term vanishes because the source mixture's mean equals the
null mean. The first detectable departure is quadratic in the shift, giving
a fourth-order likelihood divergence. This positive series also avoids
numerical cancellation in the equivalent three-term closed expression.

Including all sources and the trusted target, the unconditioned joint
chi-square divergence is
\[
 \chi_{\rm joint}^2=(1+\chi_s^2)^K
                       (1+\kappa\delta^2)^{n_t}-1.
\]
Let $q_K=\Prb\{\operatorname{Binomial}(K,w)<g\}$. Conditioning the alternative
prior to enforce the valid count changes its data law by at most $q_K$ in
total variation. Unlike the preceding broad-class construction, only the
alternative prior needs conditioning; the null is already entirely valid.

\begin{theorem}[All-valid length required by local weak-bias honesty]
\label{thm:local-lower}
For either experiment, any interval with coverage at least $1-\alpha$ at
every fixed parameter configuration in the stated local class satisfies
\[
 \E_{P_0}|I|\ge
 \delta\left[1-2\alpha-
    \min\{1,\tfrac12\sqrt{\chi_{\rm joint}^2}\}-q_K\right]_+.
\tag{L2}
\]
For fixed $\rho<w$, $M>0$, $\alpha<1/2$ and sufficiently small constant $c$,
this yields $\E_{P_0}|I|\gtrsim r_{n,K}$ as $K\to\infty$. The assertion is
at the all-valid null, not merely at an unspecified worst-case configuration.
In the i.i.d. randomized experiment it applies to honest population-TATE
intervals in this constant-effect submodel.
\end{theorem}

\begin{proof}
Uniform coverage implies coverage under the count-conditioned alternative
prior. Transfer its coverage event to $P_0$ using total variation at most
$\frac12\sqrt{\chi_{\rm joint}^2}+q_K$. Together with null coverage, both
target effects belong to the interval with probability at least the bracket
in (L2), forcing length at least $\delta$.

For the rate, set $L=\max(1,\lambda)$. Since the two-point shift multiplier
has second moment $\lambda$,
$|w+(1-w)(-\lambda)^\ell|\le\lambda L^{\ell-2}$ for $\ell\ge2$.
Using $\binom{n_s}{\ell}\le n_s^\ell/\ell!$ in (L1) gives
\[
 \chi_s^2\le\frac{\kappa^2\lambda^2}{2}
       n_s^2\delta^4\exp(\kappa n_s L^2\delta^2).
\]
Our choice of $r_{n,K}$ ensures $K n_s^2\delta^4\le c^4$,
$n_s\delta^2\le c^2K^{-1/2}$ and $n_t\delta^2\le c^2$. Therefore
$\log(1+\chi_{\rm joint}^2)=O(c^4+c^2)$ uniformly along the sequence.
Moreover $q_K\le e^{-2K(w-\rho)^2}$. A sufficiently small $c$, also satisfying
$c\le M(1-w)$, keeps the bracket bounded away from zero. Bernoulli means
remain interior for all sufficiently large sample sizes/source counts.
\end{proof}

\subsection{What this settles, and what it does not}

When $n_t\asymp n_s\asymp n$ and $K\to\infty$, the lower bound is
$n^{-1/2}K^{-1/4}$. It rules out uniformly honest oracle pooling length
$n^{-1/2}K^{-1/2}$ even at all-valid data, when nearby weakly invalid
configurations remain in the coverage class. In the exact-linear conditional
experiment, it matches the order of the main dispersion upper bound.

For the population method, fixed dimension, stable designs, constant true
CATE and $K=o(n^2)$ make its extra $n^{-1}$ composition term negligible
relative to $n^{-1/2}K^{-1/4}$. When the bad-design contribution is also
negligible, the main upper bound and this population-submodel lower bound
match in order at the null. This is a specific honest-adaptation statement,
not an optimality theorem for every CATE, design or growing dimension.

The broad bounded-invalid class contains these local linear alternatives,
so its honest procedures also face this lower bound. The result does not
show that the broader class permits a matching adaptive upper bound. In
particular, the current $n_s^{-1/2}$ bounded-invalid certificate floor remains
potentially larger than necessary at favorable configurations. Sharpening
that upper bound is a separate research problem.
